\documentclass[10pt,twocolumn]{article}
\usepackage[T1]{fontenc}
\usepackage[utf8]{inputenc}
\usepackage{lmodern}
\usepackage[margin=1.8cm,top=2.4cm,bottom=2cm,columnsep=0.6cm]{geometry}
\usepackage{fancyhdr}
\usepackage{titlesec}
\usepackage{booktabs}
\usepackage{amsmath,amssymb,amsfonts}
\usepackage{microtype}
\usepackage{xcolor}
\usepackage{mdframed}
\usepackage{parskip}
\usepackage{enumitem}
\usepackage{hyperref}

\definecolor{rulered}{RGB}{180,30,30}
\definecolor{boxgray}{RGB}{245,245,245}
\definecolor{keywordgray}{RGB}{100,100,100}

\pagestyle{fancy}
\fancyhf{}
\fancyhead[L]{\textsc{\small Claude Mag}}
\fancyhead[C]{\small\textit{Machine Learning Research Digest}}
\fancyhead[R]{\small 2026-07-27}
\fancyfoot[C]{\small\thepage}
\renewcommand{\headrulewidth}{0.8pt}

\titleformat{\section}{\normalsize\bfseries\color{rulered}}{}{0pt}{}%
  [\vspace{-4pt}\color{rulered}\rule{\columnwidth}{0.4pt}]
\titlespacing{\section}{0pt}{8pt}{4pt}

\hypersetup{colorlinks=true,urlcolor=rulered,linkcolor=black}

\begin{document}
\setlength{\parindent}{0pt}
\setlength{\parskip}{4pt}

{\small\color{keywordgray}\textbf{Keywords:}
  large-scale training \textbullet{}
  mixture-of-experts \textbullet{}
  Ascend NPU \textbullet{}
  post-training \textbullet{}
  operations research AI}\par\vspace{4pt}

{\LARGE\bfseries\raggedright
  SLAI T-Rex: Full-Parameter Post-training of the DeepSeek-V4
  Family on Ascend SuperPOD}\par\vspace{4pt}

{\large\itshape\raggedright
  Huawei Cloud Achieves a 2.93$\times$ MFU Improvement Over the
  GPU Baseline for Trillion-Parameter MoE Post-Training on Ascend
  NPUs, Demonstrating Hardware-Agnostic Frontier LLM Training}\par\vspace{6pt}

{\small\textbf{By} SLAI Team (Huawei Cloud AI Lab)}\par\vspace{2pt}

{\small\color{keywordgray}arXiv:2607.20145 \textbullet{} July 2026}
\par\vspace{2pt}\noindent{\color{rulered}\rule{\columnwidth}{0.6pt}}\vspace{6pt}

Training frontier large language models has been synonymous with NVIDIA
GPU clusters. A new technical report from Huawei Cloud challenges this
assumption: SLAI T-Rex demonstrates that full-parameter post-training
of DeepSeek-V4---a 671-billion-parameter Mixture-of-Experts model---can
be performed efficiently on Ascend NPU SuperPOD clusters via a
hierarchical optimization framework spanning model parallelism,
communication scheduling, and low-level kernel engineering. The system
achieves 34.22\% Model FLOPs Utilization (MFU)---a 2.93$\times$
improvement over the open-source GPU baseline---and produces an
operations-research specialist model that outperforms GPT-5.4-Mini.

\begin{mdframed}[backgroundcolor=boxgray,linecolor=rulered,linewidth=1pt,
    innerleftmargin=6pt,innerrightmargin=6pt,
    innertopmargin=6pt,innerbottommargin=6pt]
\textbf{\small AT A GLANCE}\par\vspace{4pt}
\begin{description}[leftmargin=0pt,labelsep=4pt,itemsep=2pt,parsep=0pt,
    font=\small\bfseries]
  \item[Problem:] Full-parameter post-training of trillion-parameter
    MoE models on non-GPU hardware suffers from memory pressure,
    communication overhead, and kernel inefficiency.
  \item[Why it matters:] Hardware diversity in large-scale LLM
    training reduces dependency on a single chip vendor and expands
    the ecosystem of frontier model development.
  \item[Method:] Expert Shard Parallelism, communication-aware
    scheduling, and Ascend-specific CANN kernel optimizations,
    combined into a three-level hierarchical framework.
  \item[Result:] 34.22\% MFU (2.93$\times$ over baseline); OR-specialist
    model achieves 71.81\% zero-shot Pass@1, +3.98\,pp over GPT-5.4-Mini.
  \item[Evidence:] Validated on DeepSeek-V4-Flash (671B active params)
    across a full CPT+SFT post-training workflow.
\end{description}
\end{mdframed}

\section{The Big Picture}

DeepSeek-V4 is a Mixture-of-Experts model with 671 billion total
parameters (37 billion active per token), trained with 4096 experts
grouped into routed and shared expert layers. Post-training this model
from scratch on Ascend hardware requires solving problems at three
intersecting layers: how to distribute the model's enormous weight
across memory-constrained NPU nodes, how to orchestrate the
expert-parallel all-to-all communications that differ fundamentally
from NVLink-based topologies, and how to write efficient NPU kernels
for the specific operations (gated attention, top-$k$ routing, grouped
GEMM) that account for the bulk of compute time.

\section{Why Existing Approaches Fall Short}

The open-source DeepSeek-V4 training recipe was designed for NVIDIA
GPU clusters with NVLink interconnects. Porting it to Ascend SuperPOD
directly yields only $\sim$11.7\% MFU because:

\begin{itemize}[itemsep=2pt,parsep=0pt]
  \item Ascend's ScaleUp network has different bandwidth/latency
    characteristics than NVLink, making the GPU pipeline schedule
    suboptimal.
  \item Expert parallelism requires all-to-all communication patterns
    that expose the Ascend cluster's inter-pod bandwidth as a bottleneck.
  \item Standard CUDA-derived kernels are mapped inefficiently to
    Ascend's cube tile architecture, leaving the NPU's SIMD units
    underutilized.
\end{itemize}

\section{Core Mechanism}

SLAI T-Rex addresses these challenges through a \textbf{three-level
hierarchical optimization framework}:

\textbf{Level 1 --- Model Parallelism:} Expert Shard Parallelism (ESP)
shards each expert's weight matrix across multiple NPU pods, reducing
per-device peak memory. A modified pipeline schedule overlaps expert
computation with inter-pod communication by reordering micro-batches
to ensure that communication and compute streams are never both idle.

\textbf{Level 2 --- Communication Orchestration:} A communication-aware
scheduler predicts expert load imbalance from recent routing statistics,
adjusts micro-batch boundaries to reduce hot-expert serialization, and
overlaps all-to-all token dispatches with local attention computations
via asynchronous execution streams.

\textbf{Level 3 --- Kernel Engineering:} Custom CANN kernels implement
fused gating (softmax + top-$k$ selection), Flash Attention adapted
to Ascend's HBM tile layout, and gradient all-reduce with dynamic
bucketing tuned to Ascend's memory bandwidth profile.

\section{Technical Formulation}

For an MoE layer with $E$ experts and top-$k$ routing, the routing
gate selects:

\[
g(x) = \mathrm{TopK}\!\left(\mathrm{softmax}(W_g x),\, k\right)
\]

and the layer output is:

\[
\mathrm{MoE}(x) = \sum_{e \in \mathrm{TopK}(x)} g_e(x)\cdot \mathrm{FFN}_e(x)
\]

Expert parallel training communicates $O(B \cdot k \cdot d)$ token
embeddings per step (where $B$ is batch size and $d$ is hidden
dimension) across all expert pods. The SLAI framework fuses the
dispatch (input routing) and combine (output aggregation) all-to-all
operations and prefetches expert weights from HBM, reducing effective
communication latency by 38\% compared to the baseline.

MFU is defined as the ratio of observed compute to peak theoretical
compute:

\[
\mathrm{MFU} = \frac{\text{Observed TFLOPs/sec}}{
  \text{Peak Theoretical TFLOPs/sec (Ascend cluster)}}
\]

The SLAI T-Rex system achieves $\mathrm{MFU} = 34.22\%$ on DeepSeek-V4-Flash
post-training, compared to $11.68\%$ for the open-source GPU recipe
adapted na\"ively to Ascend.

\section{Learning or Inference Procedure}

The OR-specialist model development follows a two-stage workflow:

\begin{enumerate}[itemsep=2pt,parsep=0pt]
  \item \textbf{Continued Pretraining (CPT):} DeepSeek-V4-Flash is
    continued-pretrained on a curated corpus of operations research
    literature, textbooks, LP/MILP problem libraries, and synthetic
    problem--solution pairs.
  \item \textbf{Supervised Fine-tuning (SFT):} 10,000 high-quality
    SFT samples spanning four task categories (LP, MILP, scheduling,
    routing) and three problem representations (mathematical, natural
    language, structured JSON) are used to align the model's output
    format and reasoning style.
\end{enumerate}

The final model is evaluated zero-shot at Pass@1 without any
inference-time search.

\section{What the Guarantee Says}

The paper does not provide theoretical guarantees on MFU or task
accuracy. The MFU improvement is empirically validated on a
full-scale DeepSeek-V4-Flash post-training run, and the 2.93$\times$
figure represents end-to-end wall-clock throughput improvement
including all communication and memory overhead.

\section{Experimental Findings}

\begin{center}
\small
\begin{tabular}{lcc}
\toprule
Model & MFU & OR Pass@1 \\
\midrule
GPU baseline (open-source) & 11.68\% & --- \\
SLAI T-Rex & 34.22\% & --- \\
\midrule
DeepSeek-V4-Flash (base) & --- & 60.54\% \\
GPT-5.4-Mini & --- & 67.83\% \\
\textbf{SLAI T-Rex OR-specialist} & --- & \textbf{71.81\%} \\
\bottomrule
\end{tabular}
\end{center}

The 2.93$\times$ MFU improvement translates directly to proportionally
reduced wall-clock training time and therefore training cost on the
same hardware budget. The OR-specialist achieves highest zero-shot
accuracy among evaluated models on a held-out OR benchmark spanning
diverse problem types and sizes.

\section{Ablations and Interpretation}

\textbf{ESP alone vs.\ full framework:} Expert Shard Parallelism
alone raises MFU from 11.68\% to 19.4\%; adding the communication
scheduler raises it to 27.1\%; adding Ascend-specific kernels
reaches 34.22\%.

\textbf{CPT vs.\ SFT-only:} The OR-specialist without CPT (SFT-only
on the base model) achieves 66.3\% Pass@1, 5.5\,pp below the full
CPT+SFT model, confirming that domain vocabulary and mathematical
notation pretraining is material.

\textbf{Data representation:} Models trained on all three problem
representations outperform single-format training by 2.1--3.8\,pp,
suggesting that diverse input format exposure improves generalization.

\section*{Reference}

SLAI Team.
\textbf{SLAI T-Rex: Full-Parameter Post-training of the DeepSeek-V4
Family on Ascend SuperPOD}.
arXiv:2607.20145, July 2026.
\url{https://arxiv.org/abs/2607.20145}

\end{document}
